%%%%%%%%%%  scan idx 223  |  folio 210  %%%%%%%%%%
est un
épimorphisme. D'autre part, on a (T.F. II 3.3.1)~:
\[
\begin{aligned}
  \Homop\bigl(i_s^{-1}(E),i_s^{-1}(F)\bigr)
  &\simeq H^0\bigl(X_s,i_s^{-1}\underline{\Homop}(E,F)\bigr)\\
  &=\varinjlim_{S'\ni s} H^0\bigl(X|S',\underline{\Homop}(E,F)\bigr)\qquad (S'\ \text{ouvert})\\
  &\simeq \varinjlim_{S'\ni s}\Homop(E|X|S',F|X|S').
\end{aligned}
\]
Donc l'isomorphisme canonique $i_s^*(E)\simeq i_s^*(F)$ est induit par un homomorphisme $f:E\to F$ (quitte à rétrécir $S$). Comme $f_s$ est un isomorphisme, $f$ est un isomorphisme au voisinage de $X_s$ par Nakayama, donc au-dessus de $X|S'$, où $S'$ est un voisinage convenable de $s$, cqfd.

\subsection*{3. Le théorème de finitude}

\subsubsection*{3.0. Conventions et notations}
\textbf{3.0.1.} Soient $X$ et $Y$ des espaces topologiques. On note $C(X,Y)$ l'ensemble des applications continues de $X$ dans $Y$, et $\underline C(X,Y)$ le faisceau sur $X$ défini par $U\mapsto C(U,Y)$.

\textbf{3.0.2.} Soient $E,F$ des espaces de Banach (réels ou complexes). On note $\Lop(E,F)$ l'espace de Banach des applications linéaires continues de $E$ dans $F$.

\textbf{3.0.3.} Les variétés mixtes dont il sera question dans ce numéro seront supposées de classe $C^\infty$. Soit $X$ une $S$-variété mixte, nous identifierons souvent un $\cO_X$ (resp. $\cO_X\otimes\mathbb C$)-Module localement libre de type fini au fibré vectoriel réel (resp. complexe) correspondant.

\textbf{3.0.4.} Soit $X$ une variété (ordinaire) compacte, et soit $E$ un fibré vectoriel (réel ou complexe) de rang fini sur $X$. On note $C^\infty(X,E)$ l'espace de Fréchet des sections (de classe $C^\infty$) sur $X$ de $E$. Pour $r\in\Z$, on note $C^r_2(X,E)$ l'espace de Hilbert des sections de $E$ $r$ fois différentiables dans $L^2$ ([16]~IX). Pour $s\ge r$, $C^s_2(X,E)$ s'envoie dans $C^r_2(X,E)$ par une application linéaire continue injective d'image dense, et l'on a 
